\section{Introduction}
\label{sec-intro}

Ontologies such as the Gene Ontology~\cite{ashburner2000} and the schemas of
enterprise knowledge graphs are stored as sets of triples and queried in
three ways.
A \emph{pattern} query asks for the triples that match a template, such as
the resources of one type, and is the base case of every RDF
store~\cite{weiss2008,neumann2010}.
A \emph{structural} query asks for the instances at or below a class, and
needs the class hierarchy.
A \emph{semantic} query asks for the terms closest in meaning to a given
term, which has no answer in the data and needs a similarity measure over
embeddings.
The three meet in queries such as ``the terms most similar to
\emph{apoptotic process} that lie under \emph{cell death}''.

Existing systems leave four problems open.
\textbf{C1. Hierarchy costs space or time.}
RDFS entailment is either \emph{materialized} at write time, which stores
every inferred triple~\cite{rdfox}, or \emph{rewritten} into the query,
which scans once per subclass.
On the Gene Ontology with 200,000 typed entities, materializing the closure
multiplies the triples by 8.1, and rewriting reads 48 times more pages than
a store that knows the hierarchy.
\textbf{C2. A sorted index does not hold meaning.}
Triple stores answer patterns with range scans over sorted permutations.
It is tempting to believe that neighbors in meaning are neighbors in the
order; we prove they cannot all be (Section~\ref{sec-order}), and measure how
few are.
Similarity is therefore left to a vector database next to the triple store,
with its own storage, its own paging, and no view of the hierarchy.
\textbf{C3. A class filter is correlated with the vectors.}
Filtered vector search is usually costed by assuming that the filter is
independent of the similarity.
In an ontology it is not: the neighbors of a term lie under the term's own
ancestors.
A post-filtered index can then find 7\% of the right answers, and the
independence estimate of the work it must do is wrong in both directions
(Section~\ref{sec-vectors}).
\textbf{C4. A paged store must scale across threads.}
A classic buffer pool serializes every page access on one latch, and a
classic B+ tree write latches the root.
In our first version, lookups per second fell from 330 thousand with one
thread to 180 thousand with eight, and inserts from 184 thousand to 23
thousand: a store embedded in a service or a notebook that answers many
queries at once spent its time waiting for itself.

\begin{example}
\label{exa-motivation}
Let $q$ be the term \emph{apoptotic process} in the Gene Ontology and $c$ the
class \emph{cell death}.
The query asks for the five terms nearest to $q$ in embedding space among
the 97 classes at or below $c$ in an 834-class slice of the ontology.
A triple store alone cannot rank by meaning.
A vector index alone returns neighbors from anywhere in the ontology, and
a post-filter keeps those under $c$ only if enough of them are.
\sys finds the 97 classes with one range scan of the same B+ tree that holds
the triples, scores them exactly, and returns \emph{apoptotic process
involved in development} (cosine 0.865) first, in under a millisecond in a
web page.
\end{example}

Database systems have handled a similar split for decades: an access method
is added to the engine that owns the storage, paging, and logging, and not
beside it.
This paper takes that view of ontologies.

\stitle{Contributions}.
This paper introduces \sys, an open-source ontology database in C++20
(Figure~\ref{fig-arch}), with five technical contributions.
\begin{enumerate}
\item \textbf{A concurrent B+ tree triple store} (Sections~\ref{sec-store}
  and~\ref{sec-conc}). Three permutations behind a buffer pool, with the
  algorithms for search, insert, and delete, a bound on the pages a range scan
  reads, and a latch-free hit path and an optimistic writer descent that scale
  to eight threads.
\item \textbf{Hierarchies as intervals} (Section~\ref{sec-labels}). A
  postorder labelling of the class DAG kept in the same B+ tree, so that
  ``everything under $c$'' is a few range scans and no entailed triple is
  stored.
\item \textbf{An entailment-aware plan rewrite} (Section~\ref{sec-rewrite}).
  RDFS subclass entailment of \texttt{rdf:type} patterns as interval scans the
  optimizer costs and joins like any other, with the answers of
  materialization, prepared 16 times faster and in 8 times less space.
\item \textbf{Vector indexes and filtered search}
  (Sections~\ref{sec-order} and~\ref{sec-vectors}). A proof that no sorted
  order keeps neighbors close, IVF posting lists stored as key ranges of the
  triple tree, and the finding that a class filter is correlated with
  similarity, with the rule that chooses between pre-filtering and masked
  search.
\item \textbf{A build for the browser} (Section~\ref{sec-system}). The whole
  engine, with its paged trees and semantic search, as a 0.85\,MB WebAssembly
  module.
\end{enumerate}
We evaluate \sys against Oxigraph~\cite{oxigraph}, Apache Jena
TDB2~\cite{jenatdb}, and SQLite on LUBM~\cite{lubm} and the Gene Ontology
(Section~\ref{sec-exp}, summarized in Table~\ref{tab-glance}).
The code, data generators, and benchmark programs are available at the
artifact URL above.
